\chapter{Hardware Design}\label{chap:hw}
The design of the robotic system focuses on accessibility, replicability, and gentle handling of the trading cards. 
This system utilizes a vacuum-based pick-and-place mechanism. This chapter details the component selection, the 
pneumatic subsystem, circuit design, and the mechanical construction.

\section{Component Selection} % TODO
The selection of components was driven by the availability of parts on the consumer market and the cost-to-performance 
ratio. A summary of key components is provided in 
\reftab{tab:components}.
\begin{table}[!h]
    \centering
    \caption{List of key hardware components.}
    \label{tab:components}
    \begin{tabular}{l|l|l}
        \textbf{Component} & \textbf{Model/Type} & \textbf{Primary Function} \\ \hline
        Microcontroller & Arduino Mega 2560 & Main control unit \\
        Stepper motor & & \\
        Stepper motor driver & TMC2209 & \\
        Rotary position sensor & AS5147-TS\_EK\_AB & Rotation feedback \\
        Servo motor & MG90s & \\
        Compressor & Sparmax TC-610H Plus & Air pressure source \\
        Venturi generator & custom 3D print & Vacuum generator \\
        Suction cup & SB 15,5 mm silicon, bellows & Card handling \\
        Pneumatic valve & Solenoid valve G356NC-18B & Controlling the airflow \\
        Noise suppressor & AE-M5 & Noise reduction \\
        Air hoses & Polyuretan, 6/4 mm & 
    \end{tabular}
\end{table}

\subsection{Control Electronics}

\subsubsection{Microcontroller: Arduino Mega}
The \textit{Arduino Mega 2560} was selected as the central processing unit. Although modern 32-bit boards (like ESP32) 
offer more power, the Mega provides a high number of GPIO pins (54 digital), which allows for easy expansion without 
multiplexing. Its 5 V logic level simplifies interfacing with certain industrial sensors, and its vast community support 
aligns with the project's educational goals.

\subsubsection{Stepper Motor and Driver}
A standard stepper motor combined with a \textit{TMC2209} driver from Trinamic is used to move the robotic arm. 
The \textit{TMC2209} driver allows configuration of motor current and micro-steps via UART, eliminating the need to work 
with a potentiometer directly on the driver. 

\subsubsection{Feedback Loop: AS5147 Encoder}
To ensure precise positioning, the system is equipped with an \textit{AS5147} magnetic rotary encoder. This sensor 
communicates via SPI and provides absolute angular position. By combining the stepper motor with this encoder, the 
system operates in a closed-loop manner, allowing it to detect missed steps or jams, which is critical for preventing 
damage to valuable cards.

\subsection{Pneumatic System}

Handling trading cards mechanically (e.g., with rubber rollers) carries a risk of scratching surfaces or bending 
corners. Therefore, a pneumatic vacuum system was designed to lift cards from the top.

\subsubsection{Air Source: Sparmax TC-610H Plus} % TODO:

\subsubsection{Vacuum Generation and Handling}
Since the compressor generates positive pressure, a \textit{Venturi generator} is used to convert compressed air into a 
vacuum. A \textit{3/2-way solenoid valve} (24 V DC, Normally Closed) controls the airflow. The 3/2 configuration is 
essential: in the \enquote{off} state, it not only stops the supply but also vents the vacuum line to the atmosphere,
allowing the card to drop instantly. Silicone suction cups with bellows are used to transfer contact with the cards 
themselves. Silicone was chosen for its softness and adhesion. Bellows were chosen to help compensate for imperfections 
in the contact between the arm and the card.

% \includesvg[inkscapeformat=pdf]{fig/sketch/magstep.svg}

\includepdf{fig/sketch/output.pdf}